% English translation of checked_pages/esq_119_025.tex; source: cote 119, site 25, PDF page 26, typescript page 23.
% AI-translated by GPT-6 Astra (OpenAI, ChatGPT Pro mode).
\scanpage{119/25}
\mspage{23}

of the 2-cartographic group that arise from regular 2-polyhedra
over finite fields (*) ? Would one obtain them all,
and, if so, how ?

These reflections bring into full light the fact,
which for me was entirely unexpected, that the theory of
finite regular polyhedra, already in the case of dimension
$n = 2$, is infinitely richer, and in particular gives infinitely
more different combinatorial forms, when one admits
base fields of nonzero characteristic,\quad than in the
case considered up to now, where the base fields were restricted
to $\mathbb{R}$ ,or, at a stretch, $\mathbb{C}$ (in the case of what Coxeter
calls ``complex regular polyhedra'', which I prefer
to call ``regular pseudo-polyhedra defined over $\mathbb{C}$ '') (**). Moreover,
it seems that this broadening of viewpoint should also cast
new light on the cases already known. Thus, examining the
Pythagorean polyhedra one after another, I saw
the same little miracle repeated each time, which I called the \uline{combinatorial
paradigm} of the polyhedron in question. Speaking vaguely, it can be
described by saying that, on looking at the specialization of the polyhedron
in the characteristic, or one of the characteristics, that is
the most singular one (or ones) (these are characteristics 2 and 5\frapparatus{V-FR-025-001}

\rule{5cm}{0.4pt}

(*) These are, moreover, the same as those arising from regular
polyhedra over any fields, or over algebraically closed fields, as
one sees by standard specialization arguments.

(**) Pinned pseudo-polyhedra are described in the same way as
pinned polyhedra, with the sole difference that the fundamental
reflections $\sigma_i$ $(0 \leq i \leq n)$ are replaced here by \uline{pseudo-reflections}
(which Coxeter moreover assumes to have finite order, as he
restricts himself to finite combinatorial structures). This simply leads
to the introduction, for each of the \unsure{.}$\sigma_i$, of an additional numerical
invariant $\beta_i$ , so that the universal n-pseudo-polyhedron can still be
defined over a polynomial ring with integer coefficients, in
the n+(n+1) variables $\underline{\alpha}_i$ $(1 \leq i \leq n)$ and $\underline{\beta}_j$ $(0 \leq j \leq n)$.

\margin{25}

